\documentclass[dvipsnames,table]{beamer}

\usetheme{Frankfurt}
\usecolortheme{orchid}

\usepackage{listings}
\usepackage{ucs}
\usepackage[utf8x]{inputenc}
\usepackage{wasysym}

\setbeamertemplate{navigation symbols}{}
\setbeamertemplate{caption}[numbered]
\setbeamerfont{caption}{size=\scriptsize}
\setbeamercolor{framenote}{bg=NetBSD-orange!25}
\setbeamercolor{rednote}{bg=Red!25}
\setbeamercolor{palette primary}{use=structure,fg=white,bg=NetBSD-orange}
\setbeamercolor{palette secondary}{use=structure,fg=white,bg=NetBSD-orange2}

\setbeamertemplate{itemize item}{\scriptsize\raise1pt\hbox{\donotcoloroutermaths$\blacktriangleright$}}
\setbeamertemplate{itemize subitem}{\tiny\raise1pt\hbox{\donotcoloroutermaths$\bullet$}}
\setbeamertemplate{itemize subsubitem}{\tiny\raise1pt\hbox{\donotcoloroutermaths{--}}}

\setbeamertemplate{enumerate item}{\insertenumlabel.}
\setbeamertemplate{enumerate subitem}{\insertenumlabel.\insertsubenumlabel}
\setbeamertemplate{enumerate subsubitem}{\insertenumlabel.\insertsubenumlabel.\insertsubsubenumlabel}
\setbeamertemplate{enumerate mini template}{\insertenumlabel}

\setbeamercolor{itemize item}{fg=NetBSD-orange, bg=NetBSD-orange}
\setbeamercolor{itemize subitem}{fg=NetBSD-orange, bg=NetBSD-orange}
\setbeamercolor{itemize subsubitem}{fg=NetBSD-orange, bg=NetBSD-orange}

\setbeamercolor{section number projected}{fg=white,bg=NetBSD-orange}
\setbeamercolor{subsection number projected}{fg=white,bg=NetBSD-orange}
\setbeamertemplate{section in toc}[circle]
\setbeamertemplate{subsection in toc}[square]

\definecolor{NetBSD-orange}{RGB}{242,103,17}
\definecolor{NetBSD-orange2}{RGB}{177,76,12}
\hypersetup{colorlinks=true,linkcolor=white,urlcolor=NetBSD-orange}

\setlength{\tabcolsep}{8pt}
\renewcommand{\arraystretch}{1.2}

\newcommand{\nbsdcolor}[1] {
	{\color{NetBSD-orange} #1}
}

\lstset{
   language=bash,
   basicstyle=\footnotesize,
   breaklines=true,
   escapechar=\@,
   commentstyle=\color{NetBSD-orange}
}

\title{Writing the NetBSD drivers\\with bus\_space(9) framework}

\author{Radoslaw Kujawa -- rkujawa@NetBSD.org}

\institute{The NetBSD Foundation}

\begin{document}

\begin{frame}
\titlepage
\end{frame}

\begin{frame}[allowframebreaks]
\frametitle{Table of Contents}
{
\hypersetup{colorlinks=true,linkcolor=black,urlcolor=NetBSD-orange}
\tableofcontents
}
\end{frame}

\section{Introduction}
\subsection{Why was this tutorial created?}

\begin{frame}
\frametitle{Why was this tutorial created?}

\begin{itemize}
	\item Introductory-level documentation is scarce
	\item Writing the device drivers is often considered black magic
	\item Reading the {\tt man} pages won't give you the big picture
	\item BSD systems are always in need of new drivers
	\item Device drivers are fun {\Large \smiley}
\end{itemize}
\end{frame}

\subsection{What won't be covered here?}

\begin{frame}
\frametitle{What won't be covered here?}

We don't have much time, so several advanced topics were omitted:

\begin{itemize}
	\item Interrupt handling
	\item Direct Memory Access and {\tt bus\_dma} framework
	\item Power management
	\item Driver detachment
	\item Drivers as kernel modules
	\item Examples for buses other than PCI
	\item Pretty much everything else...
\end{itemize}

However, once you finish this tutorial, you should be able to purse this knowledge yourself.

\end{frame}

\subsection{What is a driver anyway?}

\begin{frame}
\frametitle{What is a driver anyway?}

\begin{itemize}
	\item Interface between user space and hardware, implemented as a part of the kernel
	\item The NetBSD drivers are written mostly in C
	\item Sometimes they have machine dependent assembler parts, but this is a rare case

\end{itemize}

\end{frame}

\subsection{What do you need to write a driver?}

\begin{frame}
\frametitle{What do you need to write a driver?}

\begin{itemize}
	\item C programming skills
	\item Hardware documentation or ability to reverse engineer the hardware
	\item Reference driver implementation will help but is not essential
	\item The NetBSD installation and kernel source or cross-build environment (usually preferred for development of drivers)
	\item Some\^{}H\^{}H\^{}H\^{}H A lot of time and coffee {\Large \smiley}
\end{itemize}
\end{frame}


\section{The NetBSD driver model}

\subsection{The NetBSD kernel basics}

\begin{frame}
\frametitle{The NetBSD kernel basics}
\begin{itemize}
	\item The NetBSD has a classic monolithic kernel - all drivers are running in the same address space
	\item Thanks to above, communication between drivers and other kernel layers is simple
	\item However, it also means that one badly written driver can affect the whole kernel
	\item Numerous in-kernel frameworks standardise the way drivers are written ({\tt bus\_space}, {\tt autoconf} etc.) 
\end{itemize}
\end{frame}

\subsection{Tree shit}

\begin{frame}
\frametitle{Saving work with machine independent drivers}

\begin{itemize}
	\item If possible drivers should work on any platform
	\item High quality, machine-independent (MI) drivers are an important factor that adds to the NetBSD portability
	\item Some drivers are completely MI, some have MD or bus dependent attachments and some are completely MD
	\begin{itemize}
		\item Driver for a typical PCI card will be completely MI
		\item Driver for a SoC components will usually be completely MD
	\end{itemize}
	\item The {\tt bus\_space} abstractions helps achieving portability, transparently handling endianness issues and hiding bus implementation details from the device driver
	\item Even if we have MI drivers, writing the drivers is always significant part of effort needed to port the NetBSD to a new hardware platform
\end{itemize}
\end{frame}

\subsection{Kernel auto-configuration framework}

\subsection{Skeleton of the driver}

\section{Example driver from scratch}

\subsection{Development environment}

\begin{frame}
\frametitle{Development environment}
\begin{itemize}
	\item Out of scope of this course, but very well documented
	\item Cross compiling is an easy task with the {\tt build.sh} script
	\item Described in Part V of the NetBSD Guide
	\item Check out the NetBSD sources
	\item {\tt \$ build.sh -m cobalt tools} will build compiler, assembler, linker, etc. for cobalt port
	\item {\tt \$ build.sh -m cobalt kernel=GENERIC} will build GENERIC kernel for cobalt
	\item Pass {\tt -u} parameter to {\tt build.sh} to update (it won't rebuilding everything)
\end{itemize}
\end{frame}

\subsection{Quick introduction to {\tt GXemul}}

\begin{frame}
\frametitle{Quick introduction to {\tt GXemul}}
\begin{itemize}
	\item A framework for full-system computer architecture emulation, excellent for educational purposes
	\item Capable of emulating several real machines supported by the NetBSD
	\item We'll emulate Cobalt, a MIPS-based micro server with PCI bus
	\item I've modified {\tt GXemul} and implemented an emulation of additional PCI device
	\item It will be used to show a real-life (almost) example of driver development
\end{itemize}
\end{frame}

\subsection{Our hardware - a fake PCI card}

\begin{frame}
\frametitle{Our hardware - a fake PCI card}
\begin{itemize}
	\item Business applications often use arithmetic operations like addition
	\item FakeCards Inc. responded to market needs and created a new product, Advanced Addition Accelerator
	\item Pointy Haired Bosses will certainly buy it to accelerate their business applications, so let's create a driver for the NetBSD!

	\item Implemented as a PCI board 
\end{itemize}
\end{frame}

\subsection{Adding a new driver to the NetBSD kernel}

\subsection{Matching the PCI device}

\begin{frame}[fragile]
\frametitle{Matching the PCI device}
\begin{verbatim}
unmatched vendor 0xfabc product 0x0001 (Co-processor 
processor, revision 0x01) at pci0 dev 12 function 0 
not configured
\end{verbatim}

\begin{itemize}
	\item Kernel does not know anything about this vendor and device
	\item Add it to PCI device database
\end{itemize}
\end{frame}


\subsection{Attaching to the PCI device}

\begin{frame}
\frametitle{Attaching to the PCI device}
\begin{itemize}
	\item Our driver successfully attached but still is doing nothing useful, time to change that.
\end{itemize}
\end{frame}

\subsection{Variable types used with {\tt bus\_space}}

\begin{frame}
\frametitle{Variable types used with {\tt bus\_space}}

\begin{itemize}
         \item {\tt bus\_space\_tag\_t} -- type used to describe a particular bus, usually passed to driver from MI bus structures
         \item {\tt bus\_space\_handle\_t} --  used to describe a mapped range of bus space, usually created with {\tt bus\_space\_map()} function
         \item {\tt bus\_addr\_t} -- address on the bus
         \item {\tt bus\_size\_t} -- an amount of space on the bus
         \item Contents of these types are MD, so avoid modifying from within driver\footnote{although you'll often have to use {\tt bus\_size\_t}}
\end{itemize}
\end{frame}


\subsection{Mapping the hardware resources}

\begin{frame}[fragile]
\frametitle{Mapping the hardware resources}

\begin{verbatim}
bus_space_map(bus_space_tag_t space, bus_addr_t address, 
bus_size_t size, int flags, bus_space_handle_t  *handlep);
\end{verbatim}

\begin{itemize}
	\item The {\tt bus\_space\_map} creates mapping from physical address to virtual kernel address
         \item {\tt space} -- represents bus on which the mapping will be created
         \item {\tt address} -- typically represents physical address for which mapping will be created
	\item {\tt size} -- describes amount of bus space to be mapped
	\item {\tt handlep} -- pointer to mapped space (after successful mapping)
         \item Separate space and address
\end{itemize}
\end{frame}




\subsection{Accessing the hardware registers}
\subsection{Testing}

\section{Interacting with the user space}
\subsection{Device files}
\subsection{Utilising ioctls}
\subsection{Example user space program}



\section{A few tips}

\subsection{Avoiding common pitfalls}

\begin{frame}
\frametitle{Avoiding common pitfalls}

\begin{itemize}
	\item Always free resources allocated in {\tt match} or {\tt probe} function
	\item Always use {\tt bus\_space} methods, don't access the hardware using pointer dereference!
	\item If possible test on more than one hardware architecture, some bugs may surface
\end{itemize}
\end{frame}

\subsection{Basic driver debugging}

\begin{frame}
\frametitle{Basic driver debugging} 
\begin{itemize}
	\item Use {\tt aprint\_debug} to print debug-level messages on console and log them (enabled by passing {\tt AB\_DEBUG} from boot loader)
	\item Use built-in DDB debugger
	\begin{itemize}
			\item Enabled by kernel option {\tt DDB}
			\item Kernel panic will start DDB if {\tt DDB\_ONPANIC=1} kernel option is specified or {\tt ddb.onpanic} sysctl is set to 1.
			\item {\tt \# sysctl -w kern.panic\_now=1} to trigger panic manually ({\tt DIAGNOSTIC} option)
	\end{itemize}	
	\item Remote debugging is possible on some ports
	\begin{itemize}
			\item With KGDB through serial port
			\item With IPKDB through network
	\end{itemize}	
\end{itemize}
\end{frame}


\section{Summary}

\begin{frame}
\frametitle{Further reading}
\begin{itemize}
	\item Documentation, articles:
	\begin{itemize}
		\item A Machine-Independent DMA Framework for NetBSD, Jason R. Thorpe
		\item Writing Drivers for NetBSD, Jochen Kunz
		\item NetBSD Documentation: Writing a pseudo device
		\item driver(9), bus\_space(9), bus\_dma(9), pci(9) {\tt man} pages, etc.
	\end{itemize}
	\item Source code of the drivers:
	\begin{itemize}
		\item {\tt tdvfb}, {\tt voodoofb} are fairly good examples with documentation publicly available.
	\end{itemize}
\end{itemize}

\end{frame}

\begin{frame}
\frametitle{Questions}

\begin{itemize}
	\item Do you have any questions?
\end{itemize}
\end{frame}

\begin{frame}
\frametitle{End\ldots}
\vspace*{-0.8cm}
\begin{center}
\includegraphics[scale=0.5]{NetBSD.png}

Thank you!
\end{center}
\end{frame}

 
\end{document}
